\documentclass[12pt]{article}
\usepackage{amsmath, amssymb, amsthm} 
\usepackage[spaces, hyphens]{url}
\usepackage{hyperref}

% Customisation

\usepackage[letterpaper, top=1in, bottom=1in, left=1.2in, right=1.2in, heightrounded]{geometry}

\setlength{\parindent}{0pt}
\setlength{\parskip}{0.6em}

\title{Evaluating Randomised Algorithms and Their LLM Generated Equivalents on the Bin Packing Problem}
\author{Joshua Raybould}

\begin{document}

\maketitle

\section{Introduction}

\subsection{Problem Background}

The bin packing problem (BPP) is one of the central problems in combinatorial optimisation. The basic setup is that we are given some number of items with corresponding values (generally size or weight) and we have access to as many bins as we desire, each with the same capacity. The goal is to find a way to pack the items into the bins such that we minimise the number of bins used. This is the one-dimensional (1D) bin packing problem which is very similar to the 1D cutting stock problem (\cite{maxence}). When it comes to the one-dimensional (1D) bin packing problem (which is what we will be handling), there are two main variants. First is the offline problem, where the algorithm being used has access to all items. Second is the online problem, where 1 item is given to the algorithm at a time, at which point it must decide where to place it, meaning it has access to less information. In this instance, we will be looking at the offline case.  

This problem falls into the category NP-Hard and it would take exponential time to arrive at the exact solution. For this reason, it is a good idea to use heuristic or randomised algorithms that don't generally give the exact solution but can obtain a near-optimal one efficiently. Other than being interesting due to its innate difficulty, bin packing also has many real-world applications, such as filling up trucks that have a weight capacity and in areas like cutting industries among others (\cite{munien}). Additionally, problems in scheduling often contain it as a special case. 

\subsection{Problem Definition}
Here we are looking at the classical (offline) 1D BPP. The following is a relatively standard definition arrived at with reference to \cite{munien}, though note often an equivalent normalised version is used (see: \cite{maxence}). Each bin has the same capacity c, and each item i takes on some weight/size $s_{i}$ such that $s_{i}$ $v \leq c$ c. First, define u to be an upper bound on the minimum number of bins needed to pack all n items. Use $y_{j}$ to represent whether we have opened the jth bin, with 1 $v \leq c$ j $v \leq c$ u. Specifically, $y_j$ = 1 if bin j is open (used) in the solution otherwise it is 0. Items are labeled with i, 1 $v \leq c$ i $v \leq c$ n and we introduce $x_{ij}$ to be 1 if item i is packed in bin j. Our goal is:
\[
\min\sum_{j=1}^{u} y_j 
\]
Minimise the number of bins opened. Subject to constraints:
\[
\begin{aligned}
\sum_{j=1}^{u} x_{ij} = 1, \quad \forall i = 1, \dots, n \\
\sum_{i=1}^{n} x_{ij} s_{i}  \leq C  , \quad \forall j = 1, \dots, u
\end{aligned}
\]
Which is to say that each item must be in one and only one bin and no bin can exceed capacity.

\subsection{Randomised Algorithms}
Randomised algorithms are those that use random numbers to influence their choices, therefore they are non-deterministic. We will define stochastic optimisation as being techniques and algorithms that make use of randomness to find a near optimal solution as in \cite{sean}. It should be noted however that the more common definition requires we have a noisy objective function and/or use random choices to direct the search for a solution, see . With the definition we use we can then say the algorithms we will implement belong to the field of stochastic optimisation and in most cases a subfield of it called metaheuristics. A metaheuristic is a problem-independent strategy which guides and modifies heuristics to arrive at better solutions (). 

There are a very large number of randomised algorithms, which vary drastically in complexity.  The simplest may involve using a naive algorithm and randomising the order of the input. This is useful if you suspect the order of the input may be poor for your algorithm, for example if it was produced by an adversary. On the other hand, there are complex metaheuristics which draw inspiration from evolution and organisms in general. Part of the goal of this paper is to cover a variety of these approaches. 

\subsection{Complexity Analysis}

\subsection{Literature Review}

\subsection{Project Goals}
We plan to explore a variety of randomised algorithms that can be applied to this problem, implementing each, then analysing and comparing their performance. 

The bin packing problem has been studied extensively, and a very large number of algorithms have been proposed and applied to the problem, due to this, it is reasonable to expect that generative AI like ChatGPT will also be able to produce algorithms to solve it when prompted. This provides an opportunity to compare a solution arrived at by a human with access to the literature to that of generative AI. Comparing these will be an additional goal of this project. 

\section{Comparison Criteria}

\subsection{Criteria}
The main comparison here will be between the randomised algorithms I create, here I will define exactly how the algorithms will be compared with one another. Where possible, results will be presented using graphs.

\subsubsection{Metrics}
Generally, with algorithms, the first performance measures which come to mind are time and space complexities. These ideas are of importance and will be used where possible but as we are dealing with randomised algorithms it may be less meaningful so it will be useful to obtain empirical data too. Moreover, randomised algorithms do not arrive at exact solutions, as such, measuring how good the algorithm solutions are is arguably the most important metric. A further measure of interest is the complexity of the code. 

\subsubsection{Time}
Where possible, it makes sense to determine time complexity. I will also run each algorithm many times on the inputs to determine the mean runtime. By carrying this out for a variety of input sizes we can get a better view of each algorithms behaviour and determine when it is best applied. It may also be interesting to determine how consistent the runtime is. 

\subsubsection{Space}
Where possible, I will determine space complexity. In recent times this is not as important of a measure, so it is unlikely that we will carry out empirical testing. 

\subsubsection{Solution Quality}
Preferably, we want to compare the algorithm's solution to the best possible solution. For this, we will use ratios, as this will allow for comparison across different input lengths and make it easy to obtain a mean performance score. We will use a ratio of the number of bins used by the algorithm to the number of bins used in the optimal case. That is, if our algorithm when applied to a list of items L uses Alg bins and L when packed optimally requires Opt bins, then our ratio is Alg/Opt. It is clear from this that the greater the ratio the worse the algorithm has performed. We can then take the mean of this ratio over a variety of cases to give an idea of how it tends to perform. This is limited since if an algorithm performs well in just a certain case, it may be missed or skew the results. Due to this, I will also take performance measures for specific sizes and plot them and/or determine the mean Alg/Opt for ranges of sizes.

\subsubsection{Code Complexity}
This measure will aim to look at how hard it was to implement the algorithm and how elegant the code is. The number of lines of code used will be taken into account, but more important is how long it took to implement, if it required parameters to be tuned and how understandable the code is just by looking at it. 

\subsection{Comparison Of AI generated algorithms}
As an additional goal, I will compare my algorithms to the same kind of algorithm generated by ChatGPT's GPT-5. This comparison will be carried out in the same way as outlined above using the same inputs, but I will also briefly compare the code to see if either was simpler and reflect on what it has produced to see if I can learn something from it. By incorporating this as an additional goal it is important to set boundaries on how the AI is going to being used. For starters, the model we will be using will be GTP-5. The prompts will generally consist of asking for a specific algorithm to solve the problem in Python, then if it does not work, prompting it again with information about the issue until either it arrives at a solution or fails to do so. In the case it fails to do so, if possible, I will then fix it and have this noted as a caveat of the AI's solution. Additionally, whilst I will not be using AI for the research or implementation of my own solutions, there are currently no limitations on how else I use it, for example I may use it to help in deciding which algorithm to implement next. If I feel these decisions affect the comparisons in some way, they will be made known. 

\subsection{The Problem Instances}
To ensure this is as comprehensive as possible, a wide variety of inputs will be used. It makes sense to use existing instances since they have solutions allowing us to compare the algorithm's solution to the optimal solution, it also means the performance of algorithms here is more likely to be able to be compared to those in other studies (as they may use the same data). The instances will be taken from BPPLIB (see:\cite{BPPLIB}). Specifically, I will use the Randomly Generated instances as these cover a range of input lengths, item weights and bin capacities. Additionally, I will use Falkenauer's instances, in particular the set of 80 instances which uses uniformly distributed item sizes and also ANI AI as they are difficult instances. This will be supplemented with instances that I randomly generate, which will have the option to be either uniformly or normally distributed. Here, the minimum possible weight of an item will be between 10\% and 30\% of the capacity of the bins (with the percentage chosen at random) and the maximum possible weight be equivalent to the capacity of the bins. 

\section{The Algorithms}

\subsection{Best Fit With Randomised Order}
\subsubsection{Overview}
Best Fit is a very typical solution for the online version of the bin packing problem but isn't usually applied to the offline version (since there are better alternatives). It makes sense to start with this, as it is one of the simplest algorithms that could be applied to the problem that will still serve as a good baseline of performance. This will be the normal algorithm where the item is placed in the smallest available space it can fit into but with the order of the input items being randomised. The purpose of this is to allow the algorithm to (generally) do better when given bad orderings \cite{rajeev}.
\subsubsection{Time and Space Complexity}
In all cases randomising the input using the Fisher-Yates shuffle requires O(n) operations. The best fit algorithm then runs. In the worst case each item is too large to fit in any open bin meaning the algorithm will on average search through n/2 bins before opening a new bin. This happens for all n items, hence O($n^2$). Therefore, the overall worst case time complexity is O($n^2$).

As for space complexity (ignoring the space for the solution since all algorithms will require approximately the same) it is O(1) as it just requires 2 variables to carry out the randomisation. 

\subsection{Simulated Annealing}
\subsubsection{Overview}
This is extremely similar to hill-climbing, which works like so:

Start with some candidate solution.

Repeat until we arrive at a local optimum or run out of time:
\begin{enumerate}
\item Alter the candidate solution slightly (move to a neighbouring solution).
\item If this altered solution is better than our current candidate solution, then 
select it to be our new candidate solution; otherwise, there is no change.
\end{enumerate}

The difference with simulated annealing is that rather than only selecting a new solution if it is better than the current candidate solution, it can with some probability choose to move to a worse solution. The purpose of this is to provide an opportunity to escape a local optimum, increasing the chance of finding and converging at the global optimum. This probability depends partially on a temperature value; the higher this is the more likely the algorithm is to move to worse solutions. The temperature decreases as the algorithm progresses on a schedule we decide. When it comes to researching algorithms, my main sources are (TODO cite the books, maybe move this cite).

\subsubsection{Implementation}
Firstly, the initial candidate solution is generated by using a first fit decreasing (FFD) algorithm as in \cite{sonuc}. We also use the same objective function and their ideas for moving from a candidate solution to a neighbouring one (note these are quite common ideas, see for example: \cite{kryz}). Specifically, the objective function will be:
\[
\max\sum_{j=1}^{u} \sum_{i=1}^{n} {(x_{ij} s_{i})}^{2}
\]
In the actual implementation this is divided by bin capacity. Note that we can improve this function without decreasing the number of bins used, but it effectively directs the algorithm towards solutions where the bins are as full as possible. To move to a neighbouring solution, we select two different bins (a and b) at random. Then we decide whether we will move an item from bin a to bin b or if we will swap items between them. There will be an equal chance of either happening and the items to move will also be selected at random. If the move/swap would lead the solution to violate a constraint then the candidate solution remains unchanged.

Moreover, we need to define the general parameters present in simulated annealing. We set temperature = 1000000 and decide on using an exponential cooling schedule with our  $\alpha$ = 0.995. This schedule was chosen as it was more effective than a linear cooling schedule, likely due to its slower cooling allowing for more exploration (see \cite{algafter} for more about cooling schedules). This value for $\alpha$ is large, but using one lower than this meant it was uncommon for solutions to actually be improved in terms of bins due FFD being relatively strong (see \cite{dosa}  where it is shown that the solution FFD comes to is at worst $\tfrac{11}{9}$ times the optimal plus $\tfrac{6}{9}$).

\subsubsection{Time and Space Complexity}

Space complexity is O(1)

\subsection{Grouping Genetic Algorithm}

\subsubsection{Overview}

The genetic algorithm (GA) is a population-based method and takes ideas from genetics and evolution, so classes as an Evolutionary Algorithm (EA). Generally, the GA takes a desired population size and generates a population of that size consisting of randomly generated individuals (candidate solutions). The ideal population size is problem dependent. If it is too small, then there will not be enough variety in the population to explore the search space sufficiently, but if it is too large, it will run too slowly. Next, the algorithm repeatedly selects an individual from the population based on some strategy and applies crossover and/or mutation to it. The specifics on this vary, it can be guaranteed that both crossover and mutation happen or each happen with some probability etc. The selection method is most often tournament selection - which effectively takes a tournament size t as input and selects the best of t randomly chosen individuals from the population based on fitness (note that the selection does not remove the individual from the population). If our individuals are bit strings, then crossover could be carried out by swapping some bits between two individuals (parents). Mutating an individual would flip bits within it with some probability. Once the selected individual has been altered, it is added to a new population. Upon reaching the set population size, the new population is used in the process of building another new population. This continues until we reach a termination condition, which is often after a certain amount of time. At which point we pick the best solution from the population. (TODO: We may want to mention elitism but this depends on if the GGA is going to use it) 

Unfortunately, the BPP is a grouping problem which GAs don't tend to perform well on, so it requires some modifications after which it should become a competitive algorithm for the problem. Specifically, we will be implementing a grouping genetic algorithm (GGA).
\section{Comparison}

\begin{thebibliography}{00}
\bibitem{maxence} Maxence Delorme, Manuel Iori, Silvano Martello, Bin packing and cutting stock problems: Mathematical models and exact algorithms, European Journal of Operational Research, Volume 255, Issue 1, 2016, Pages 1-20, ISSN 0377-2217, \url{https://doi.org/10.1016/j.ejor.2016.04.030}.

\bibitem{munien} Munien, Chanaleä and Ezugwu, Absalom E.. "Metaheuristic algorithms for one-dimensional bin-packing problems: A survey of recent advances and applications" Journal of Intelligent Systems, vol. 30, no. 1, 2021, pp. 636-663. \url{https://doi.org/10.1515/jisys-2020-0117}

\bibitem{sean} Sean Luke, 2013, Essentials of Metaheuristics, Lulu, second edition, available for free at \url{http://cs.gmu.edu/~sean/book/metaheuristics/ }

\bibitem{brownlee} Brownlee, J. (2021, October 12). A gentle introduction to stochastic optimization algorithms. Machine Learning Mastery. \url{https://machinelearningmastery.com/stochastic-optimization-for-machine-learning/}

\bibitem{spall} J.~C.~Spall, ``Stochastic Optimization,'' in \emph{Handbook of Computational Statistics: Concepts and Methods}, 2nd ed., J.~Gentle, W.~Härdle, and Y.~Mori, Eds. Heidelberg: Springer-Verlag, 2012, ch.~7, pp.~173--201. \url{https://doi.org/10.1007/978-3-642-21551-3_7}

\bibitem{kamali} Kamali, Shahin, and Pooya Nikbakht. "Cutting Stock with Rotation: Packing Square Items into Square Bins." International Conference on Combinatorial Optimization and Applications. Springer, Cham, 2020. \url{https://mspace.lib.umanitoba.ca/items/1d1c973b-1dbb-46c3-bb80-c07fe746db35}. Kamali, Shahin, and Pooya Nikbakht. "On the Fault-Tolerant Online Bin Packing Problem." arXiv preprint arXiv:2107.02922 (2021). 

\bibitem{BPPLIB} M. Delorme, M. Iori, and S. Martello. BPPLIB: A library for bin packing and cutting stock problems. Optimization Letters, 12(2):235-250, 2018. \url{https://github.com/mdelorme2/BPPLIB}.

\bibitem{rajeev}

\bibitem{sonuc} Sonuç, E., Şen, B. and Bayır, Ş., 2017. Solving Bin Packing Problem Using Simulated Annealing. In: Proceedings of the 65th ISERD International Conference, Mecca, Saudi Arabia, 23–24 January 2017. World Research Library, pp. 83–85. \url{https://worldresearchlibrary.org/up\_proc/pdf/635-148739346283-85.pdf}

\bibitem{kryz} Krzysztof Fleszar, Khalil S. Hindi, New heuristics for one-dimensional bin-packing, Computers \& Operations Research, Volume 29, Issue 7, 2002, Pages 821-839, ISSN 0305-0548, \url{https://doi.org/10.1016/S0305-0548(00)00082-4}.

\bibitem{algafter} Algorithm Afternoon, n.d. Chapter 2: Temperature and Cooling Schedules. In: Simulated Annealing Afternoon: A Practical Guide for Software Developers. AlgorithmAfternoon.com. Available at: \url{https://algorithmafternoon.com/books/simulated\_annealing/chapter02/}

\bibitem{dosa} D\'osa, G. (2007). The tight bound of First Fit Decreasing bin-packing algorithm is FFD(I) $\leq$ 11/9 OPT(I) + 6/9. In B. Chen, M. Paterson, \& G. Zhang (Eds.), \textit{Combinatorics, Algorithms, Probabilistic and Experimental Methodologies (ESCAPE 2007)} (Lecture Notes in Computer Science, vol.~4614, pp.~1--11). Springer, Berlin, Heidelberg. \url{https://doi.org/10.1007/978-3-540-74450-4\_1}



\end{thebibliography}
\end{document}
